\chapter{Cơ sở lý thuyết: Nhận diện biển số xe}
\label{ch:ly_thuyet_alpr}

\section{Tổng quan về ALPR}
Hệ thống nhận diện biển số xe tự động (Automatic License Plate Recognition - ALPR) thường bao gồm hai giai đoạn chính:
\begin{enumerate}
    \item \textbf{Plate Detection:} Xác định vị trí biển số trong ảnh
    \item \textbf{Character Recognition (OCR):} Nhận dạng các ký tự trên biển số
\end{enumerate}

\begin{figure}[!htbp]
    \centering
    \begin{tikzpicture}[
        node distance=1cm,
        box/.style={rectangle, draw, rounded corners, minimum width=2.2cm, minimum height=0.9cm, align=center, fill=blue!10, font=\small},
        arrow/.style={->, thick, >=stealth}
    ]
        \node[box] (input) {Ảnh\\đầu vào};
        \node[box, right=of input] (detect) {Plate\\Detection};
        \node[box, right=of detect] (preprocess) {Tiền xử lý\\(CLAHE)};
        \node[box, right=of preprocess] (ocr) {OCR};
        \node[box, right=of ocr] (output) {Biển số};

        \draw[arrow] (input) -- (detect);
        \draw[arrow] (detect) -- (preprocess);
        \draw[arrow] (preprocess) -- (ocr);
        \draw[arrow] (ocr) -- (output);
    \end{tikzpicture}
    \caption{Pipeline tổng quan của hệ thống ALPR}
    \label{fig:alpr_pipeline}
\end{figure}

\section{Biển số xe Việt Nam}

\subsection{Các định dạng biển số}
Biển số xe Việt Nam có các định dạng chính sau:

\begin{table}[!ht]
\centering
\caption{Các định dạng biển số xe Việt Nam}
\label{tab:vn_plate_format}
\begin{tabular}{|c|l|l|}
\hline
\textbf{Loại} & \textbf{Định dạng} & \textbf{Ví dụ} \\ \hline
Xe ô tô (1 dòng) & XX-Y XXXXX & 30G-07784 \\ \hline
Xe ô tô (2 dòng) & XX-Y / XXXXX & 30G / 07784 \\ \hline
Xe máy & XX-YY / XXX.XX & 29-H1 / 234.56 \\ \hline
\end{tabular}
\end{table}

Trong đó:
\begin{itemize}
    \item 2 chữ số đầu: Mã tỉnh/thành phố (VD: 30 = Hà Nội, 29 = Hà Nội cũ, 51 = TP.HCM)
    \item Chữ cái: Ký hiệu seri (A-Z, trừ I, O, Q)
    \item 5 chữ số cuối: Số thứ tự đăng ký
\end{itemize}

\subsection{Thách thức nhận diện biển số Việt Nam}
Việc nhận diện biển số Việt Nam gặp một số thách thức đặc thù:
\begin{itemize}
    \item \textbf{Đa dạng định dạng:} Biển số có thể là 1 dòng hoặc 2 dòng, với các kích thước khác nhau
    \item \textbf{Font chữ đặc thù:} Font chữ trên biển số Việt Nam có một số điểm khác biệt
    \item \textbf{Điều kiện thực tế:} Biển số thường bị bẩn, mờ, hoặc nghiêng do góc quay camera
    \item \textbf{Ánh sáng:} Độ phản chiếu của biển số thay đổi theo điều kiện ánh sáng
\end{itemize}

\section{PaddleOCR - Nhận dạng ký tự}

\subsection{Giới thiệu PaddleOCR}
PaddleOCR là một bộ công cụ OCR mã nguồn mở được phát triển bởi PaddlePaddle (Baidu). PaddleOCR hỗ trợ nhận dạng văn bản đa ngôn ngữ với độ chính xác cao và tốc độ xử lý nhanh, phù hợp cho các ứng dụng thời gian thực.

\subsection{Kiến trúc CRNN}
PaddleOCR sử dụng kiến trúc CRNN (Convolutional Recurrent Neural Network) bao gồm:
\begin{enumerate}
    \item \textbf{CNN Backbone:} Trích xuất đặc trưng không gian từ ảnh đầu vào. Backbone thường sử dụng là MobileNet hoặc ResNet biến thể.
    \item \textbf{Sequence Modeling (LSTM):} Mô hình hóa quan hệ tuần tự giữa các ký tự. Sử dụng Bidirectional LSTM để capture thông tin từ cả hai hướng.
    \item \textbf{CTC Decoder:} Giải mã chuỗi ký tự cuối cùng từ output của LSTM.
\end{enumerate}

\begin{figure}[!htbp]
    \centering
    \begin{tikzpicture}[
        node distance=1cm,
        box/.style={rectangle, draw, rounded corners, minimum width=2.2cm, minimum height=0.9cm, align=center, fill=green!10, font=\small},
        arrow/.style={->, thick, >=stealth}
    ]
        \node[box] (input) {Ảnh biển số\\$H \times W$};
        \node[box, right=of input] (cnn) {CNN\\Backbone};
        \node[box, right=of cnn] (lstm) {Bi-LSTM\\$T \times D$};
        \node[box, right=of lstm] (ctc) {CTC\\Decoder};
        \node[box, right=of ctc] (output) {30G07784};

        \draw[arrow] (input) -- (cnn);
        \draw[arrow] (cnn) -- (lstm);
        \draw[arrow] (lstm) -- (ctc);
        \draw[arrow] (ctc) -- (output);
    \end{tikzpicture}
    \caption{Kiến trúc CRNN cho OCR trong PaddleOCR}
    \label{fig:crnn_arch}
\end{figure}

\subsection{CTC Loss}
CTC (Connectionist Temporal Classification) cho phép huấn luyện mà không cần alignment giữa input và output. Đây là đặc điểm quan trọng vì ta không biết chính xác vị trí của từng ký tự trên ảnh.

Xác suất của một chuỗi output $\mathbf{l}$ được tính:
\begin{equation}
    P(\mathbf{l} | \mathbf{x}) = \sum_{\pi \in \mathcal{B}^{-1}(\mathbf{l})} P(\pi | \mathbf{x})
\end{equation}

Trong đó $\mathcal{B}^{-1}(\mathbf{l})$ là tập tất cả các path có thể tạo ra $\mathbf{l}$ sau khi loại bỏ ký tự blank và ký tự lặp.

\textbf{Ví dụ:} Với chuỗi mục tiêu ``AB'', các path hợp lệ có thể là:
\begin{itemize}
    \item ``A-A-B-B'' $\rightarrow$ ``AB'' (loại bỏ lặp)
    \item ``A--B-'' $\rightarrow$ ``AB'' (loại bỏ blank `-')
    \item ``-A-B-'' $\rightarrow$ ``AB''
\end{itemize}

CTC Loss được tính bằng negative log-likelihood:
\begin{equation}
    \mathcal{L}_{CTC} = -\log P(\mathbf{l} | \mathbf{x})
\end{equation}

\section{Tiền xử lý ảnh biển số}

Chất lượng ảnh biển số ảnh hưởng trực tiếp đến độ chính xác của OCR. Chúng tôi áp dụng các kỹ thuật tiền xử lý sau:

\subsection{CLAHE - Cân bằng histogram thích ứng}
CLAHE (Contrast Limited Adaptive Histogram Equalization) là kỹ thuật cải thiện độ tương phản ảnh, đặc biệt hiệu quả cho ảnh có độ sáng không đều.

\subsubsection{Nguyên lý hoạt động}
\begin{enumerate}
    \item \textbf{Chia tile:} Chia ảnh thành các tile nhỏ (VD: 8×8 pixels)
    \item \textbf{Histogram Equalization cục bộ:} Thực hiện histogram equalization trên từng tile
    \item \textbf{Clip Limit:} Giới hạn độ tương phản (clip limit) để tránh khuếch đại noise. Nếu histogram bin vượt quá clip limit, phần dư sẽ được phân bố lại đều cho các bin khác.
    \item \textbf{Bilinear Interpolation:} Nội suy bilinear để ghép các tile lại một cách mượt mà, tránh hiện tượng blocking artifact.
\end{enumerate}

\subsubsection{Công thức toán học}
Công thức histogram equalization cho một pixel:
\begin{equation}
    s_k = (L-1) \cdot \sum_{j=0}^{k} p_r(r_j) = (L-1) \cdot \text{CDF}(r_k)
\end{equation}

Trong đó:
\begin{itemize}
    \item $s_k$: Mức xám mới của pixel
    \item $L$: Số mức xám (256 cho ảnh 8-bit)
    \item $p_r(r_j)$: Xác suất xuất hiện của mức xám $r_j$
    \item CDF: Cumulative Distribution Function của histogram
\end{itemize}

\textbf{Clip Limit trong CLAHE:}
\begin{equation}
    \beta = \frac{M \times N}{L} \times \text{clipLimit}
\end{equation}
Trong đó $M \times N$ là kích thước tile và clipLimit là tham số (thường từ 2.0 đến 4.0).

\subsection{Sharpening - Làm sắc nét}
Bộ lọc làm sắc nét giúp tăng cường các cạnh của ký tự, làm cho chúng rõ ràng hơn trước khi đưa vào OCR.

\subsubsection{Kernel Sharpening}
Chúng tôi sử dụng kernel convolution để làm sắc nét:
\begin{equation}
    K_{sharpen} = \begin{bmatrix}
        0 & -1 & 0 \\
        -1 & 5 & -1 \\
        0 & -1 & 0
    \end{bmatrix}
\end{equation}

\subsubsection{Nguyên lý}
Kernel này hoạt động bằng cách:
\begin{itemize}
    \item Nhân pixel trung tâm với hệ số lớn (5)
    \item Trừ đi trung bình của 4 pixel lân cận
    \item Kết quả: tăng cường sự khác biệt giữa pixel và các pixel xung quanh $\rightarrow$ cạnh sắc nét hơn
\end{itemize}

Công thức tính pixel mới:
\begin{equation}
    I'(x,y) = 5 \cdot I(x,y) - [I(x-1,y) + I(x+1,y) + I(x,y-1) + I(x,y+1)]
\end{equation}

\subsection{Quy trình tiền xử lý hoàn chỉnh}
\begin{enumerate}
    \item \textbf{Crop ảnh biển số:} Từ kết quả detection của YOLOv8
    \item \textbf{Chuyển Grayscale:} Để áp dụng CLAHE
    \item \textbf{Áp dụng CLAHE:} Với \texttt{clipLimit=2.0}, \texttt{tileGridSize=(8,8)}
    \item \textbf{Chuyển về RGB:} Để đưa vào PaddleOCR
    \item \textbf{Sharpening (tùy chọn):} Áp dụng nếu ảnh còn mờ
\end{enumerate}

\section{So sánh các phiên bản YOLOv8}

Trong đề tài này, chúng tôi sử dụng \textbf{YOLOv8n} (nano) cho cả phát hiện phương tiện và phát hiện biển số vì yêu cầu cân bằng giữa tốc độ và độ chính xác.

\begin{table}[!ht]
\centering
\caption{So sánh các phiên bản YOLOv8}
\label{tab:yolov8_versions}
\begin{tabular}{|c|c|c|c|}
\hline
\textbf{Model} & \textbf{Params (M)} & \textbf{FLOPs (G)} & \textbf{mAP@50 (COCO)} \\ \hline
YOLOv8n & 3.2 & 8.7 & 37.3 \\ \hline
YOLOv8s & 11.2 & 28.6 & 44.9 \\ \hline
YOLOv8m & 25.9 & 78.9 & 50.2 \\ \hline
YOLOv8l & 43.7 & 165.2 & 52.9 \\ \hline
YOLOv8x & 68.2 & 257.8 & 53.9 \\ \hline
\end{tabular}
\end{table}

\textbf{Lý do chọn YOLOv8n:}
\begin{itemize}
    \item Số lượng tham số nhỏ (3.2M) $\rightarrow$ phù hợp triển khai trên phần cứng phổ thông
    \item FLOPs thấp (8.7G) $\rightarrow$ tốc độ inference cao
    \item Đối với bài toán phát hiện biển số (single class, đối tượng tương đối dễ nhận diện), YOLOv8n đủ để đạt độ chính xác cao
\end{itemize}

\section{Cơ chế Validation \& Retry}

\subsection{Vấn đề với OCR đơn frame}
OCR có thể cho kết quả khác nhau giữa các frame do:
\begin{itemize}
    \item Biển số bị rung, blur khi xe di chuyển
    \item Ánh sáng thay đổi (phản chiếu, bóng)
    \item Góc nhìn thay đổi khi xe tiến lại gần camera
    \item Nhiễu từ mô hình OCR
\end{itemize}

\subsection{Giải pháp: Format Validation}
Thay vì lưu trữ nhiều kết quả và voting, chúng tôi sử dụng phương pháp validate format biển số Việt Nam:

\begin{algorithm}[!ht]
\SetAlgoLined
\KwIn{$ocr\_text$: Kết quả OCR, $vehicle\_class$: Loại xe}
\KwOut{$(is\_valid, cleaned\_plate)$: Kết quả validation}

$cleaned \leftarrow$ RemoveSpaces($ocr\_text$)\;
\eIf{$vehicle\_class \in \{motorbike, bicycle\}$}{
    $pattern \leftarrow$ ``\textasciicircum{}[0-9]\{2\}[A-Z][A-Z0-9][0-9]\{4,5\}\$''\;
    $min\_len \leftarrow 8$, $max\_len \leftarrow 9$\;
}{
    $pattern \leftarrow$ ``\textasciicircum{}[0-9]\{2\}[A-Z][0-9]\{4,5\}\$''\;
    $min\_len \leftarrow 7$, $max\_len \leftarrow 8$\;
}
\If{$min\_len \leq \text{len}(cleaned) \leq max\_len$ \textbf{AND} Match($pattern$, $cleaned$)}{
    \Return $(True, cleaned)$\;
}
\Return $(False, cleaned)$\;
\caption{Thuật toán Validation biển số Việt Nam}
\label{alg:validation}
\end{algorithm}

\subsection{Retry Mechanism}
Khi kết quả OCR không đúng format, hệ thống sẽ retry ở frame tiếp theo:
\begin{itemize}
    \item Chỉ lưu kết quả \textbf{đã validate} (đúng format)
    \item Nếu OCR sai format hoặc không đọc được $\rightarrow$ không lưu, thử lại frame sau
    \item Khi đã có kết quả hợp lệ $\rightarrow$ không xử lý OCR nữa (tiết kiệm tài nguyên)
\end{itemize}

\begin{equation}
    \text{ShouldRunOCR}(track\_id) = \begin{cases}
        \text{True} & \text{nếu chưa có kết quả valid} \\
        \text{False} & \text{nếu đã có kết quả valid}
    \end{cases}
\end{equation}

\section{Character Correction cho biển số Việt Nam}

\subsection{Các lỗi phổ biến của OCR}
Do một số ký tự có hình dáng tương tự, PaddleOCR có thể nhầm lẫn:

\begin{table}[!ht]
\centering
\caption{Bảng hiệu chỉnh ký tự phổ biến}
\label{tab:char_correction}
\begin{tabular}{|c|c|l|}
\hline
\textbf{Sai} & \textbf{Đúng} & \textbf{Điều kiện} \\ \hline
0 (số không) & O (chữ O) & Vị trí chữ cái (index 2) \\ \hline
O (chữ O) & 0 (số không) & Vị trí chữ số \\ \hline
1 (số một) & I (chữ I) & Vị trí chữ cái \\ \hline
8 (số tám) & B (chữ B) & Vị trí chữ cái \\ \hline
5 (số năm) & S (chữ S) & Vị trí chữ cái \\ \hline
\end{tabular}
\end{table}

\subsection{Logic hiệu chỉnh dựa trên định dạng}
Dựa trên cấu trúc biển số Việt Nam (ví dụ: 30G-07784):
\begin{itemize}
    \item \textbf{Vị trí 0-1:} Bắt buộc là chữ số (mã tỉnh)
    \item \textbf{Vị trí 2:} Bắt buộc là chữ cái (seri)
    \item \textbf{Vị trí 3-7:} Bắt buộc là chữ số (số thứ tự)
\end{itemize}

Thuật toán hiệu chỉnh:
\begin{enumerate}
    \item Kiểm tra độ dài chuỗi OCR
    \item Tại mỗi vị trí, kiểm tra xem ký tự có đúng loại (chữ/số) không
    \item Nếu sai loại, tra bảng mapping để sửa
\end{enumerate}

\section{Hai phương pháp triển khai ALPR}

Trong hệ thống của chúng tôi, có hai cách tiếp cận để thực hiện nhận diện biển số xe:

\subsection{Phương pháp 1: Two-Stage ALPR (Hai giai đoạn)}

Đây là phương pháp truyền thống, chia quá trình thành hai bước độc lập:

\subsubsection{Kiến trúc}
\begin{enumerate}
    \item \textbf{Giai đoạn 1 - Phát hiện phương tiện:} Sử dụng YOLOv8n (5 lớp) để phát hiện các loại xe (xe máy, xe đạp, ô tô, xe buýt, xe tải)
    \item \textbf{Giai đoạn 2 - Phát hiện biển số:} Crop vùng xe từ giai đoạn 1, chạy qua YOLOv8n chuyên biệt (1 lớp: license\_plate) để phát hiện biển số
    \item \textbf{Giai đoạn 3 - OCR:} Crop vùng biển số, áp dụng tiền xử lý (CLAHE, sharpening), sau đó sử dụng PaddleOCR để nhận dạng ký tự
\end{enumerate}

\subsubsection{Ưu điểm}
\begin{itemize}
    \item \textbf{Độ chính xác cao:} Mô hình chuyên biệt cho biển số đạt mAP@50 = 99.5\%
    \item \textbf{Phát hiện biển nhỏ tốt:} Vì chỉ focus vào 1 task (plate detection), mô hình học rất tốt các đặc trưng của biển số
    \item \textbf{Linh hoạt:} Có thể thay thế từng module độc lập (ví dụ: thay YOLOv8 bằng Faster R-CNN cho plate detection)
\end{itemize}

\subsubsection{Nhược điểm}
\begin{itemize}
    \item \textbf{Chậm:} Cần 2 lần inference qua YOLO (vehicle detection + plate detection), tốc độ chỉ đạt 5--10 FPS
    \item \textbf{Tiêu tốn tài nguyên:} Phải load 2 models YOLOv8 vào GPU/RAM (~2.5GB)
    \item \textbf{Phức tạp:} Logic phải quản lý 2 models, crop vùng xe, rồi mới detect plate
\end{itemize}

\subsection{Phương pháp 2: One-Stage Semi-supervised (Một giai đoạn)}

Phương pháp này tích hợp việc phát hiện phương tiện và biển số vào một mô hình duy nhất.

\subsubsection{Kiến trúc}
\begin{enumerate}
    \item \textbf{Giai đoạn 1 - Phát hiện đa lớp:} Sử dụng YOLOv8n (6 lớp) để phát hiện đồng thời cả xe và biển số trong một lần forward pass
    \item \textbf{Giai đoạn 2 - Matching:} Sử dụng thuật toán containment để gán biển số cho xe (kiểm tra xem bbox biển có nằm trong bbox xe không)
    \item \textbf{Giai đoạn 3 - OCR:} Tương tự Phương pháp 1, crop vùng biển số và nhận dạng ký tự
\end{enumerate}

\subsubsection{Quy trình gán nhãn bán giám sát (Semi-supervised)}
Để có được mô hình 6 lớp, chúng tôi áp dụng kỹ thuật Semi-supervised Labeling:
\begin{enumerate}
    \item \textbf{Chuẩn bị Teacher Model:} Sử dụng mô hình phát hiện biển số từ Phương pháp 1 (mAP@50 = 99.5\%) làm teacher
    \item \textbf{Tự động gán nhãn:} Chạy teacher model trên toàn bộ dataset vietnam\_vehicle\_dataset (chỉ có 5 lớp xe) để tự động phát hiện và gán nhãn biển số
    \item \textbf{Hợp nhất dataset:} Merge các annotation biển số vào file gốc, tạo thành dataset 6 lớp (5 loại xe + license\_plate)
    \item \textbf{Huấn luyện lại:} Train YOLOv8n từ đầu trên dataset 6 lớp này
\end{enumerate}

\subsubsection{Ưu điểm}
\begin{itemize}
    \item \textbf{Nhanh:} Chỉ cần 1 lần inference qua YOLO, tốc độ đạt 15--20 FPS
    \item \textbf{Tiết kiệm tài nguyên:} Chỉ cần load 1 model (~1.2GB)
    \item \textbf{Mapping tự động:} Dễ dàng gán biển số cho xe thông qua kiểm tra containment
    \item \textbf{Học ngữ cảnh:} Mô hình học được mối quan hệ không gian giữa xe và biển số
\end{itemize}

\subsubsection{Nhược điểm}
\begin{itemize}
    \item \textbf{Độ chính xác thấp hơn:} mAP@50 cho tất cả lớp chỉ đạt 78.5\% (do class imbalance)
    \item \textbf{Khó phát hiện biển nhỏ:} Biển số chiếm tỷ lệ nhỏ trong dataset, dễ bị miss khi xe ở xa
    \item \textbf{Phụ thuộc teacher model:} Chất lượng annotation tự động phụ thuộc vào độ chính xác của teacher model
\end{itemize}

\subsection{So sánh hai phương pháp}

\begin{table}[!ht]
\centering
\caption{So sánh hai phương pháp triển khai ALPR}
\label{tab:alpr_method_comparison_theory}
\begin{tabular}{|l|p{5cm}|p{5cm}|}
\hline
\textbf{Tiêu chí} & \textbf{Phương pháp 1 (Two-Stage)} & \textbf{Phương pháp 2 (One-Stage)} \\ \hline
Số model & 2 (Vehicle + Plate) & 1 (6-class) \\ \hline
Tốc độ & Chậm (5--10 FPS) & Nhanh (15--20 FPS) \\ \hline
Độ chính xác plate & Cao (99.5\%) & Trung bình (78.5\%) \\ \hline
Tài nguyên & Nhiều (~2.5GB) & Ít (~1.2GB) \\ \hline
Triển khai & Phức tạp & Đơn giản \\ \hline
Sử dụng tốt nhất & Yêu cầu độ chính xác cao, camera tốt & Ưu tiên tốc độ, real-time \\ \hline
\end{tabular}
\end{table}

Trong thực tế, chúng tôi triển khai cả hai phương pháp và cho phép người dùng lựa chọn phù hợp với nhu cầu của hệ thống (độ chính xác cao vs. tốc độ nhanh).
